\section{CM tensors and four-factor relations}\label{sec:tensors}

This section reduces Theorem~\ref{thm:cm} to the algebraicity of a
specified family of tensors in the fourth cohomology of products of four
CM abelian varieties.  The reduction changes the CM types of individual
factors while preserving their total Hodge type; its intermediate
factors may therefore be absent from the original variety.  We first
record how to work with individual embedding lines and recover
rationality after extending coefficients.  We then construct the
auxiliary CM varieties and assemble arbitrary balanced tensors from
four-factor relations by algebraic correspondences.

\subsection{Scalar coefficients and degree-one maps}

For a smooth projective complex variety \(X\), write
\[
 \Hdg^{2p}(X)=H^{2p}(X,\Q)\cap H^{p,p}(X),\qquad
 \mathcal Z^{2p}(X)=
 \im\bigl(\CH^p(X)_{\Q}\longrightarrow H^{2p}(X,\Q)\bigr).
\]
Cohomology is Betti cohomology unless another theory is specified.  We
fix \(\Qbar\subset\C\), and write \(T_k=T\otimes_{\Q}k\) for a rational
vector space \(T\).  A class with coefficients in \(\Qbar\) or \(\C\)
is called \emph{algebraic} if it belongs to the corresponding scalar
extension of \(\mathcal Z^{2p}(X)\).  Scalar conjugation therefore
preserves algebraicity on the same complex variety: it fixes each
rational cycle class and changes only its coefficient.  We use the
compatibility of cycle classes with exterior products, cup products,
pullback between smooth varieties, and proper pushforward.

Our Tate convention is \(V(r)^{a,b}=V^{a+r,b+r}\), so \(V(r)\) has
weight \(\operatorname{wt}(V)-2r\).

\begin{lemma}[Scalar descent]\label{lem:trace-descent}
Let \(T\) be a finite-dimensional rational vector space and
\(P\subset T_{\C}\) a complex-linear subspace.  If \(t\in T_{\Qbar}\)
and \(\sigma t\in P\) for every
\(\sigma\in\Gal(\Qbar/\Q)\), then
\[
 t\in (T\cap P)\otimes_{\Q}\Qbar.
\]
The same assertion applies to the subspace of maps of any fixed Hodge
type in a Hom space.
\end{lemma}

\begin{proof}
Choose a finite Galois extension \(K/\Q\) containing the coefficients
of \(t\).  For \(a\in K\), the vector
\[
 T_a=\Tr_{K/\Q}(at)
     =\sum_{\sigma\in\Gal(K/\Q)}\sigma(a)\sigma(t)
\]
is rational.  Each automorphism of \(K\) extends to \(\Qbar\), so every
summand belongs to \(P\); hence \(T_a\in T\cap P\).  If
\(b_1,\ldots,b_r\) is a rational basis of \(K\) and
\(b_1^*,\ldots,b_r^*\) is its trace-dual basis, write
\(t=\sum_jt_jb_j\) with \(t_j\in T\).  Then \(T_{b_j^*}=t_j\), which
proves the assertion.  The condition that a complex-linear map have a
fixed Hodge type is complex-linear, so the same argument applies in a
Hom space.
\end{proof}

We will also use
\begin{equation}\label{eq:cycle-descent}
 \bigl(\mathcal Z^{2p}(X)\otimes_{\Q}\Qbar\bigr)
       \cap H^{2p}(X,\Q)=\mathcal Z^{2p}(X).
\end{equation}
Indeed, extend a rational basis of \(\mathcal Z^{2p}(X)\) to one of
\(H^{2p}(X,\Q)\) and compare coefficients.

\begin{lemma}[Realization in degree one]\label{lem:weight-one-maps}
For complex abelian varieties \(A\) and \(B\), pullback induces an
isomorphism
\[
 \Hom(A,B)\otimes_{\Z}\Q
 \ \simeq\
 \Hom_{\mathrm{Hdg}}\bigl(H^1(B,\Q),H^1(A,\Q)\bigr).
\]
If \(S\) is a smooth connected projective complex variety and
\(h:H^1(B,\Q)\to H^1(S,\Q)\) is a rational Hodge map, then there are
an integer \(m\geq1\) and a morphism \(f:S\to B\) such that
\(f^*=mh\) on \(H^1(B,\Q)\).
\end{lemma}

\begin{proof}
Dualizing a rational Hodge map gives a rational map on homology
commuting with the complex structures.  Multiplication by a positive
integer makes it carry the integral lattice of \(A\) into that of
\(B\).  It consequently induces a holomorphic homomorphism of complex
tori \(A^{\mathrm{an}}\to B^{\mathrm{an}}\), which is algebraic.
Conversely, an algebraic homomorphism induces the indicated Hodge
pullback.  This proves the first assertion; equivalently it is the
dual of the rational form of Riemann's theorem
\cite[Proposition~6.5, Theorems~6.7--6.8, and Corollary~6.9]{MilneShimura2017}.

The Albanese map of \(S\) induces an isomorphism
\(H^1(\Alb S,\Q)\simeq H^1(S,\Q)\).  Apply the first assertion to the
map with target \(H^1(\Alb S,\Q)\), clear a denominator, and compose
the resulting homomorphism \(\Alb S\to B\) with an Albanese map.
\end{proof}

The Lefschetz \((1,1)\) theorem says that
\(\Hdg^2(X)=\mathcal Z^2(X)\).  In conjunction with
Lemma~\ref{lem:trace-descent}, it has the following useful consequence:
if \(t\in H^2(X,\Qbar)\) is of type \((1,1)\) in every scalar
conjugate, then \(t\) is algebraic.  This observation will supply all
two-factor tensors used below.

\subsection{Rank-one CM structures}

For a polarizable rational Hodge structure \(H\) of weight one, let
\(J(H)\) be the smallest rational algebraic subgroup of \(\GL(H)\)
containing the Hodge circle.  The circle acts by \(z\) on \(H^{1,0}\)
and by \(\bar z\) on \(H^{0,1}\), for \(|z|=1\).  It preserves an
alternating polarization.  Its connectedness shows that it lies in the
rational identity component of \(J(H)\); minimality therefore makes
\(J(H)\) connected.

\begin{lemma}\label{lem:toric-cm}
For a complex abelian variety \(A\), the following conditions are
equivalent:
\begin{enumerate}[label=\textup{(\roman*)}]
\item \(J(H^1(A,\Q))\) is a torus.
\item \(H^1(A,\Q)\) is a direct sum of rational Hodge structures,
      each of rank one over a CM field acting by Hodge endomorphisms.
\item \(\End(A)\otimes_{\Z}\Q\) contains a commutative semisimple
      algebra of rational dimension \(2\dim A\).
\end{enumerate}
\end{lemma}

\begin{proof}
The zero-dimensional case is immediate.  Put \(H=H^1(A,\Q)\), and
first assume that \(J(H)\) is a torus.  Its associative envelope
\(E_0\subset\End_{\Q}(H)\) is a commutative semisimple algebra: after
splitting the torus, its elements act diagonally by scalars on the
weight spaces, so the envelope is a subalgebra of a product of fields
and is reduced.  This description descends to \(\Q\).

Let \(*\) be the polarization adjoint.  It preserves \(E_0\), since
it sends every symplectic torus element to its inverse.  Every element
of \(E_0\) commutes with the complex structure \(I\) supplied by the
Hodge circle.  Consequently \(*\) is also the adjoint for the positive
definite real form obtained by inserting \(I\) in the polarization.
In particular,
\[
 \Tr_{H_{\R}}(a^*a)>0\qquad(0\ne a\in E_0\otimes_{\Q}\R).
\]
The involution cannot exchange two primitive factors of the real
finite étale algebra \(E_0\otimes\R\): for the idempotent \(e\) of one
factor, this would give \(e^*e=0\).  On a complex factor it cannot be
the identity, since for its idempotent \(e\) one would have
\((ie)^*(ie)=-e\), of negative trace.  Thus it acts as complex
conjugation on each complex factor.  There are no real factors.  On
the nonzero module of such a factor the circle would act by real
scalars, and a continuous homomorphism from the circle to
\(\R^\times\) is trivial; this contradicts the weight-one characters.

It follows that each rational field factor \(K\) of \(E_0\) is CM.
Indeed, it has only complex archimedean factors, and its involution
satisfies \(\tau(a^*)=\overline{\tau(a)}\) at every embedding.  Its
fixed field is therefore totally real and has index two.  On the
corresponding \(K\)-module the circle acts by \(K\otimes\R\)-scalars.
Splitting that module into \(K\)-lines thus gives rational Hodge
substructures.  This proves (ii), without discarding multiplicities.

The separate scalar actions on the field lines in (ii) form a
commutative semisimple Hodge-endomorphism algebra of dimension
\(\dim_{\Q}H\).  Lemma~\ref{lem:weight-one-maps} realizes it in
\(\End(A)\otimes\Q\), proving (iii).  Conversely, write an algebra as
in (iii) as \(\prod_jK_j\).  Its faithful action on \(H\) has a
positive rank \(m_j\) over each \(K_j\), and
\[
 \dim_{\Q}H=\sum_jm_j[K_j:\Q]=\sum_j[K_j:\Q].
\]
Hence every \(m_j=1\).  The centralizer of this action in \(\GL(H)\)
is the torus \(\prod_j\Res_{K_j/\Q}\mathbf G_m\).  It contains the
Hodge circle because the algebra acts by Hodge endomorphisms, proving
(i).
\end{proof}

These equivalences give the rank-one description of the CM varieties
defined in the introduction.
A finite product of CM abelian varieties is CM: the circle on the
direct sum of their degree-one cohomologies lies in the product of
their tori, and its connected rational envelope is a subtorus.

Choose a Galois CM field \(E\subset\Qbar\), and put
\[
 F=E^+,\qquad G_E=\Gal(E/\Q),\qquad s=[F:\Q].
\]
Its complex conjugation \(c\in G_E\) is central.  One such \(E\) can
contain any finite collection of CM fields and all their conjugates.
To see this, take the compositum of all those conjugates.  For a CM
field \(K\), every embedding \(\rho:K\to\C\) satisfies
\(\rho c_K=c_{\C}\rho\).  These identities show that complex
conjugation on the Galois compositum commutes with its Galois group.
Every conjugate of an element fixed by it is real, so the fixed field
is totally real.  The compositum is therefore a Galois CM field.

An \emph{odd sign function} on \(G_E\) is a function
\(v:G_E\to\{1,-1\}\) satisfying \(v(c\tau)=-v(\tau)\).  Give the
rational vector space \(E\) a weight-one Hodge structure \(U(v)\) by
declaring its embedding line
\[
 U(v)_\lambda=
 \{u\in U(v)_{\Qbar}: au=\lambda(a)u\text{ for every }a\in E\}
 \qquad(\lambda\in G_E)
\]
to have type \((1,0)\) if \(v(\lambda)=1\), and type \((0,1)\)
otherwise.  Oddness makes these declarations compatible with the real
structure.

This Hodge structure is polarizable.  Choose \(0\ne\delta\in E\) with
\(c(\delta)=-\delta\).  Multiplying \(\delta\) by an element of
\(F^\times\), weak approximation gives an element \(\xi\in E^\times\)
such that \(c(\xi)=-\xi\) and
\[
 \operatorname{sgn}\bigl(\operatorname{Im}\tau(\xi)\bigr)=-v(\tau)
\]
at one representative \(\tau\) of each conjugate pair.  The form
\[
 \psi_v(x,y)=\Tr_{E/\Q}\bigl(\xi x c(y)\bigr)
\]
is alternating and nondegenerate.  It pairs only conjugate embedding
lines.  On a vector \(u\) in the sum of the \((1,0)\) lines,
\(i\psi_v(u,\bar u)\) is a sum of positive multiples of the squared
absolute values of its embedding coordinates; hence it is positive
for \(u\ne0\).  This is the weight-one polarization condition.
After scaling, the form is integral on a lattice.  The Riemann
criterion, applied to the dual homological Hodge structure, gives a
CM abelian variety \(B(v)\) of dimension \(s\) with
\begin{equation}\label{eq:cm-types}
 U(v)=H^1(B(v),\Q).
\end{equation}
This construction permits every odd sign function, including
imprimitive CM types; it does not require a principal polarization.
The polarization and Riemann-criterion argument also applies to any
CM field, without a Galois hypothesis, and any CM type, that is, any
choice of one embedding from each conjugate pair.  The Galois field
\(E\) provides a common indexing set for the types used below.

We will need the exact action of scalar conjugation on labels.  If
\(\sigma\in\Gal(\Qbar/\Q)\) restricts to \(g\in G_E\) and
\(u\in U(v)_\lambda\), rationality of the \(E\)-action gives
\[
 a\,\sigma(u)=\sigma(au)
             =\sigma(\lambda(a))\sigma(u)
             =(g\lambda)(a)\sigma(u).
\]
Consequently
\begin{equation}\label{eq:cm-label-conjugation}
 \sigma U(v)_\lambda=U(v)_{g\lambda}.
\end{equation}
In particular the complete type signs of \(U(v)_1\) and \(U(v)_c\)
after \(g\) are \(v(g)\) and \(-v(g)\), respectively.  Here and below
the commutation of \(c\) with \(g\) takes place in \(G_E\); no
centrality assertion is made in \(\Gal(\Qbar/\Q)\).

The same convention applies to a line in a rank-one module over a CM
subfield of \(E\).  Choose \(\lambda\in G_E\) extending the embedding
that labels the line.  Its \emph{complete type function} takes the value
\(+1\) at \(g\in G_E\) if the line labeled by the restriction of
\(g\lambda\) has type \((1,0)\), and \(-1\) if it has type \((0,1)\).
It is odd.  Since the embedding-line decomposition is defined over
\(E\), this is the type of the scalar-conjugate line for every
automorphism of \(\Qbar\) restricting to \(g\).

\begin{lemma}[Transfer of an embedding line]\label{lem:cm-line-transfer}
Let \(H\) and \(H'\) be direct sums of rank-one weight-one Hodge
modules over CM fields contained in \(E\).  Choose an embedding line
\(L\subset H_{\Qbar}\) and an embedding line \(M\subset H'_{\Qbar}\)
with the same complete type function.  Extend any
\(\Qbar\)-linear map \(T:L\to M\) by zero on all other source
embedding lines.  Then
\[
 T\in\Hom_{\mathrm{Hdg}}(H,H')\otimes_{\Q}\Qbar.
\]
When \(H\) and \(H'\) are degree-one cohomologies of abelian varieties,
this map is therefore an algebraic-coefficient linear combination of
pullbacks of abelian morphisms, up to rational denominators.
\end{lemma}

\begin{proof}
Scalar conjugation transports the entire source and target
decompositions.  For each \(\sigma\), the only possibly nonzero
component of \(\sigma T\) maps \(\sigma L\) to \(\sigma M\), whose
types agree by hypothesis.  All other components are zero.  Thus
\(\sigma T\) is a Hodge map on the whole source.  Apply
Lemma~\ref{lem:trace-descent} in \(\Hom_{\Q}(H,H')\), followed, when
applicable, by Lemma~\ref{lem:weight-one-maps}.
\end{proof}

The map \(T\) in this lemma need not be invariant under a stabilizer
of either line.  This matters when a field is a proper subfield of
\(E\), or when a rational field module occurs more than once.  The
proof treats every scalar conjugate of the zero-extended map and
retains the repeated rational summands separately.

\subsection{Reduction to four factors}

The four-factor assertion that remains to be proved is the following.

\begin{theorem}[Four-factor switch]\label{thm:switch}
Let \(E\) be a Galois CM field and let
\(v_1,v_2,v_3,v_4:G_E\to\{1,-1\}\) be odd sign functions satisfying
\(v_1+v_2=v_3+v_4\).  Under the Künneth inclusion, the following line is algebraic:
\begin{equation}\label{eq:switch-line}
 U(v_1)_1\otimes U(v_2)_1\otimes U(v_3)_c\otimes U(v_4)_c
 \ \subset\
 H^4\!\left(\prod_{i=1}^4B(v_i),\Qbar\right)
\end{equation}
\end{theorem}

This line is already of Hodge type \((2,2)\) in every scalar
conjugate.  Indeed, after an automorphism restricting to \(g\in G_E\),
its four type signs have sum
\begin{equation}\label{eq:switch-conjugate-balance}
 v_1(g)+v_2(g)-v_3(g)-v_4(g)=0.
\end{equation}
The theorem asks for algebraicity, which does not follow from this
balance alone.  We leave its proof to the subsequent construction;
the present section proves that it is sufficient.

Hazama's coordinate-arrangement calculation rationally spans the
kernel of its Hodge matrix by antipodal pairs and elementary
four-label relations \cite[Theorem~7.13]{Hazama2003CM}.  His
Corollary~7.14.1 states a corresponding codimension-two algebraicity
reduction.  We prove the particular tensor reduction needed here
directly, including the transfers to auxiliary abelian factors and
the algebraic contractions that assemble the relations.

\begin{proposition}\label{prop:switch-reduction}
If Theorem~\ref{thm:switch} holds for every Galois CM field and every
quadruple of odd sign functions satisfying its relation, then
Theorem~\ref{thm:cm} holds.
\end{proposition}

\begin{proof}
Let \(A\) be a CM abelian variety and let \(h\in\Hdg^{2p}(A)\).
The cases \(p=0\) and \(p=1\) follow from the unit class and the
Lefschetz \((1,1)\) theorem, respectively.  For higher codimension,
we first replace a Hodge monomial by a tensor on auxiliary CM
factors.  We then connect its two lists of types by four-factor
switches and compose the resulting algebraic tensors.

By Lemma~\ref{lem:toric-cm}, choose
a rational decomposition of \(H=H^1(A,\Q)\) into rank-one CM-field
modules, retaining separate copies of repeated summands.  Choose a
Galois CM field \(E\) containing all the fields and their conjugates.
Over \(E\), \(H\) is the direct sum of its embedding lines.  Since cup
product gives \(H^*(A,\Q)=\Lambda^*H\), the exterior monomials in
these lines give a direct-sum decomposition of \(H^*(A,\Qbar)\) into
one-dimensional lines \cite[Proposition~6.4]{MilneShimura2017}.
Every such line is contained in one Hodge bidegree.

Scalar conjugation permutes the monomial
lines, with only the possible sign from reordering exterior factors.
Because \(h\) is rational, a nonzero component on one monomial line
gives a nonzero component on every line in its scalar-Galois orbit.
All these lines must have type \((p,p)\).  The decomposition is
defined over \(E\), so the components of \(h\) have algebraic Betti
coefficients.

Fix a supported monomial and list the complete type functions of its
\(2p\) factors as \(z_1,\ldots,z_{2p}\).  Each \(z_i\) is odd.  In the
conjugate indexed by \(g\), the number of \((1,0)\) factors minus the
number of \((0,1)\) factors is \(\sum_i z_i(g)\).  Its balance is
therefore exactly
\[
 \sum_{i=1}^{2p}z_i=0.
\]
For \(1\leq i\leq p\), set \(x_i=z_i\) and \(y_i=-z_{p+i}\).  Then
\(\sum_i x_i=\sum_i y_i\).  The first \(p\) factors have the complete
types of \(U(x_i)_1\), and the last \(p\) have those of \(U(y_i)_c\).
For any ordered list \(a=(a_1,\ldots,a_p)\) of odd sign functions,
put \(A_a=\prod_{i=1}^pB(a_i)\); for two such lists put
\[
 \mathscr L_{a,b}=
 \bigotimes_{i=1}^pU(a_i)_1\ \otimes\
 \bigotimes_{i=1}^pU(b_i)_c
 \subset H^{2p}(A_a\times A_b,\Qbar).
\]
Repeated factors in these products are kept as separate slots.
Lemma~\ref{lem:cm-line-transfer} reduces the chosen monomial to the
algebraicity of \(\mathscr L_{x,y}\).  We return to \(A\) after
composing the transition tensors.

Regard each list as a matrix with \(p\) rows and one column for each
representative of \(G_E/\langle c\rangle\).  At any column, equality
of the two sums is equality of the numbers of plus signs.  Among the
rows where the current \(x\)-entry differs from the desired
\(y\)-entry, there are equally many mismatches of the forms
\((+1,-1)\) and \((-1,+1)\).  Pair these two sets of rows and swap the
two current entries in every pair, making the opposite swap at the
conjugate coordinate.  This corrects the column, leaves the other
columns unchanged, and preserves oddness of every row.
Successively correcting the columns gives a finite sequence
\[
 x=x^{(0)},\ x^{(1)},\ \ldots,\ x^{(N)}=y
\]
in which consecutive lists differ by one such two-row swap.

At a single swap the old rows \(a,b\) and new rows \(a',b'\) satisfy
\(a+b=a'+b'\) at every coordinate.  Theorem~\ref{thm:switch} supplies
the pure algebraic tensor
\[
 U(a)_1\otimes U(b)_1\otimes U(a')_c\otimes U(b')_c.
\]
Every unchanged row is joined to its new copy by
\(U(v)_1\otimes U(v)_c\), an all-conjugate \((1,1)\) line and hence an
algebraic two-factor tensor.  Pull these classes to their distinct
old and new slots and multiply.  This gives a nonzero algebraic
generator of \(\mathscr L_{x^{(j)},x^{(j+1)}}\).  If the original
lists already agree, the two-factor tensors alone generate
\(\mathscr L_{x,x}\).  It remains to compose consecutive transition
tensors without losing their algebraicity or nonvanishing.

For three ordered lists \(a,b,e\), suppose that
\(\mathscr L_{a,b}\) and \(\mathscr L_{b,e}\) have pure algebraic
generators with algebraic Betti coefficients.  Write them as
\[
 \begin{split}
 t_{ab}&=\alpha_1\otimes\cdots\otimes\alpha_p\otimes
          \beta_1\otimes\cdots\otimes\beta_p,\\
 t_{be}&=\gamma_1\otimes\cdots\otimes\gamma_p\otimes
          \delta_1\otimes\cdots\otimes\delta_p,
 \end{split}
\]
where \(\beta_i\in U(b_i)_c\) and \(\gamma_i\in U(b_i)_1\).  Identify the
middle copies and choose an ample divisor class \(\ell_i\) on each
\(B(b_i)\).  With \(r_{ab},r_{be},r_{ae}\) denoting the projections
from \(A_a\times A_b\times A_e\), and with the \(\ell_i\) also denoting
their pullbacks, cup product and pushforward give
\begin{equation}\label{eq:transition-composition}
 \begin{split}
 (r_{ae})_*\left(
 r_{ab}^*t_{ab}\smile r_{be}^*t_{be}
       \smile\prod_{i=1}^p\ell_i^{\,s-1}\right)
 &=
 (-1)^{p(p-1)/2}\left(\prod_{i=1}^p q_i\right)t_{ae},\\
 q_i&=\int_{B(b_i)}\beta_i\wedge\gamma_i\wedge\ell_i^{\,s-1},
 \end{split}
\end{equation}
where \(t_{ae}=\alpha_1\otimes\cdots\otimes\alpha_p\otimes
\delta_1\otimes\cdots\otimes\delta_p\).  The sign is obtained by regrouping
\(\beta_1\cdots\beta_p\gamma_1\cdots\gamma_p\) into
\(\beta_1\gamma_1\cdots\beta_p\gamma_p\) in the indicated slot order.

Each \(q_i\) is nonzero.  Complex conjugation exchanges the two
middle embedding lines, so \(\gamma_i\) is a nonzero scalar multiple of
\(\bar\beta_i\).  Hodge--Riemann positivity for the ample class makes
the pairing with \(\bar\beta_i\) nonzero, whether \(\beta_i\) is holomorphic
or antiholomorphic.  This requires no compatibility between \(\ell_i\)
and \(E\).  Moreover \(q_i\in\Qbar^\times\), since its vectors have
algebraic Betti coefficients and \(\ell_i\) is rational.  The
codimension on the left of Equation~\eqref{eq:transition-composition}
is
\[
 p+p+p(s-1)-ps=p.
\]
Thus it is an algebraic class of the required degree, and division by
the nonzero algebraic coefficient shows that the endpoint line is
algebraic.  Iterating along the sequence of lists proves algebraicity
of \(\mathscr L_{x,y}\).

Finally choose the line transfers in the direction from the auxiliary
\(U\)'s to \(H^1(A,\Q)\).  Each is an algebraic-coefficient linear
combination of rational pullbacks along maps from \(A\) to the
corresponding \(B(v)\).  Apply their exterior product on \(A^{2p}\)
and pull back along the diagonal \(A\to A^{2p}\).  The result is the
cup product of the selected original factors.  They occupy distinct
embedding or multiplicity lines in the original exterior monomial,
so this cup product is nonzero.  We have proved that every supported
monomial line is algebraic.  Summing the components of \(h\) and
using Equation~\eqref{eq:cycle-descent} proves the rational assertion.

Finite products are covered because they are themselves CM,
as observed after Lemma~\ref{lem:toric-cm}.
\end{proof}

The proof uses tensors on auxiliary products and contracts their
intermediate slots before returning to \(A\).  It does not assert
generation of the Hodge ring on a fixed power by cup products of
classes in codimensions one and two.

\begin{lemma}[Choice of the distinguished embedding]
 \label{lem:switch-normalization}
It suffices to prove Theorem~\ref{thm:switch} when \(s\geq2\) and
\(v_i(1)=1\) for all four functions.
\end{lemma}

\begin{proof}
If \(v_1\) and \(v_2\) are nowhere simultaneously positive, oddness
also rules out simultaneous negativity.  Hence \(v_2=-v_1\).
The relation gives \(v_4=-v_3\), and the first two and the last two
slots of Equation~\eqref{eq:switch-line} are separately
all-conjugate \((1,1)\) pairs.

Otherwise choose \(g\in G_E\) at which \(v_1(g)=v_2(g)=1\).  The
relation makes all four signs positive there.  Choose a scalar
automorphism extending \(g\).  It sends the labels in
Equation~\eqref{eq:switch-line} to \(g,g,gc,gc\).  The complete type
function of the line labeled \(g\) is the right translate
\[
 v_i^g(\tau)=v_i(\tau g).
\]
It is odd, satisfies \(v_i^g(1)=1\), and the four translated functions
satisfy the same relation.  Centrality of \(c\) in \(G_E\) makes the
complete type of the \(gc\)-line equal to \(-v_i^g\).  The
transfers of Lemma~\ref{lem:cm-line-transfer} therefore carry the
normalized switch for the \(v_i^g\)'s to the lines labeled
\(g,g,gc,gc\) on the original factors.  Applying the inverse scalar
automorphism gives Equation~\eqref{eq:switch-line}.

When \(s=1\), there is only one odd sign function positive at the
identity, so the normalized tensor is a product of ordinary pairs.
\end{proof}

For the remaining construction we work in the normalized case, so
all four lines \(U(v_i)_1\) are holomorphic.  The next section
isolates a surface-period condition sufficient to prove the switch.
